\documentclass[a4paper,11pt]{article}

\newcommand{\TPnumero}{Project 5}
\newcommand{\macrosPath}{../../preamble/macros.tex}
% =========================================================
% Tutorial / lab preamble — Time Series Analysis module
% article class + fancyhdr + tcolorbox + listings (English)
% Author : Aymen Ben Brik (aymen.benbrik@esprit.tn)
% =========================================================

\usepackage[utf8]{inputenc}
\usepackage[T1]{fontenc}
\usepackage[english]{babel}
\usepackage{lmodern}
\usepackage{amsmath, amssymb, amsthm, mathtools, bm}
\usepackage{graphicx}
\usepackage{booktabs}
\usepackage{array}
\usepackage{multirow}
\usepackage{colortbl}
\usepackage{xcolor}
\usepackage{tikz}
\usetikzlibrary{positioning, arrows.meta, calc, shapes, fit, backgrounds, matrix}
\usepackage{listings}
\usepackage[most]{tcolorbox}
\usepackage{geometry}
\geometry{a4paper, margin=2cm, headheight=15pt}
\usepackage{fancyhdr}
\usepackage[hidelinks]{hyperref}
\usepackage{enumitem}

% --- Colours ---
\definecolor{espritBlue}{RGB}{30, 60, 110}
\definecolor{espritAccent}{RGB}{200, 80, 30}
\definecolor{boxEx}{RGB}{30, 60, 110}
\definecolor{boxQ}{RGB}{180, 120, 0}
\definecolor{boxC}{RGB}{30, 100, 60}
\definecolor{boxAtt}{RGB}{200, 80, 30}

% --- Listings ---
\definecolor{codebg}{RGB}{248,248,248}
\definecolor{codeframe}{RGB}{220,220,220}
\definecolor{keyword}{RGB}{0,82,155}
\definecolor{comment}{RGB}{88,110,117}
\definecolor{string}{RGB}{133,153,0}

\lstdefinestyle{pythonstyle}{
  language=Python,
  basicstyle=\ttfamily\small,
  keywordstyle=\color{keyword}\bfseries,
  commentstyle=\color{comment}\itshape,
  stringstyle=\color{string},
  backgroundcolor=\color{codebg},
  frame=single,
  rulecolor=\color{codeframe},
  frameround=tttt,
  breaklines=true,
  showstringspaces=false,
  tabsize=2
}
\lstdefinestyle{Rstyle}{
  language=R,
  basicstyle=\ttfamily\small,
  keywordstyle=\color{keyword}\bfseries,
  commentstyle=\color{comment}\itshape,
  stringstyle=\color{string},
  backgroundcolor=\color{codebg},
  frame=single,
  rulecolor=\color{codeframe},
  frameround=tttt,
  breaklines=true,
  showstringspaces=false,
  tabsize=2
}
\lstset{style=pythonstyle}

% --- Common macros (configurable path) ---
\providecommand{\macrosPath}{../../preamble/macros.tex}
\input{\macrosPath}

% --- Exercise counter ---
\newcounter{excounter}
\renewcommand{\theexcounter}{\arabic{excounter}}

\newtcolorbox[use counter=excounter]{exercise}[1][]{
  enhanced, breakable,
  colback=boxEx!5, colframe=boxEx, fonttitle=\bfseries,
  title={Exercise~\theexcounter\ifstrempty{#1}{}{ -- #1}},
  arc=2pt, boxrule=0.6pt
}
\newtcolorbox{question}[1][]{
  enhanced, breakable,
  colback=boxQ!5, colframe=boxQ, fonttitle=\bfseries,
  title={Question\ifstrempty{#1}{}{ -- #1}},
  arc=2pt, boxrule=0.5pt
}
\newtcolorbox{solution}[1][]{
  enhanced, breakable,
  colback=boxC!5, colframe=boxC, fonttitle=\bfseries,
  title={Solution\ifstrempty{#1}{}{ -- #1}},
  arc=2pt, boxrule=0.5pt
}
\newtcolorbox{warning}[1][]{
  enhanced, breakable,
  colback=boxAtt!5, colframe=boxAtt, fonttitle=\bfseries,
  title={Warning\ifstrempty{#1}{}{ -- #1}},
  arc=2pt, boxrule=0.5pt
}

% --- Headers / footers ---
\pagestyle{fancy}
\fancyhf{}
\fancyhead[L]{\textsc{\moduleNom} -- \auteurNom}
\fancyhead[R]{\TPnumero}
\fancyfoot[C]{\thepage}
\fancyfoot[L]{\auteurInstitution}
\fancyfoot[R]{\auteurEmail}
\renewcommand{\headrulewidth}{0.4pt}
\renewcommand{\footrulewidth}{0.2pt}

% --- Placeholder ---
\providecommand{\TPnumero}{Lab}


\title{Project 5: Full Box--Jenkins ARMA study with forecasting\\
\large Identification $\cdot$ Estimation $\cdot$ Diagnostics $\cdot$ Forecast}
\author{\auteurNom\\\auteurEmail\\\auteurInstitution}
\date{\anneeUniversitaire}

\begin{document}
\maketitle

\section*{Context}

This project consolidates all of Chapters~3--5: stationarisation,
ACF/PACF identification, ARMA estimation, residual diagnostics and
forecasting. You will produce a complete Box--Jenkins study on the
Souvenir Sales data, on the differenced log-scale, where an
ARMA model is appropriate.

The deliverable serves two purposes: it will be re-used in
Chapter~6 (ARCH/GARCH) where you will check the residuals for
volatility clustering.

\section*{Deliverables}
\begin{itemize}
  \item Reproducible \texttt{project5\_<name>.py} with seed $42$.
  \item Report PDF (8--10 pages) including all figures, all
        numerical outputs, and a final back-transformed forecast
        with $95\%$ CI.
  \item Optional 10-min oral defence.
\end{itemize}

% =========================================================
\begin{exercise}[Stationarisation]
\textbf{(1)} Load \texttt{data/SouvenirSales.csv}. Compute
$y_t = \log(\text{sales}_t)$.

\textbf{(2)} Apply the seasonal log-difference
$w_t = (1 - L)(1 - L^{12}) y_t$. Verify by ADF/KPSS that $w_t$ is
stationary.

\textbf{(3)} From now on, model $w_t$ as an \emph{ARMA} on the
stationary scale. Plot $w_t$, its sample ACF and PACF up to
lag~$36$.
\end{exercise}

% =========================================================
\begin{exercise}[Identification]
\textbf{(1)} Read the ACF/PACF of $w_t$ at the regular lags
($1$--$11$): which (small) AR or MA orders look plausible?

\textbf{(2)} Read at the seasonal lags ($12, 24, 36$): does the
spike at lag~$12$ persist or disappear after one differencing?

\textbf{(3)} Propose three candidate orders, e.g.\
$(0, 1)$, $(1, 0)$, $(1, 1)$ on the regular part (and you may add a
seasonal component if needed; see Chapter~4).
\end{exercise}

% =========================================================
\begin{exercise}[Estimation and selection]
\textbf{(1)} Fit each candidate by MLE
(\texttt{statsmodels.tsa.arima.ARIMA}). For each, report:
\begin{itemize}
  \item the estimated coefficients with standard errors;
  \item the log-likelihood, AIC, BIC, $\widehat\sigma^2$;
  \item the Ljung--Box $p$-value at $m = 12$ and $m = 24$.
\end{itemize}

\textbf{(2)} Choose the best model on the (AIC, BIC, Ljung--Box)
triplet. Justify in 5 lines.

\textbf{(3)} Plot the model residuals: time plot, sample ACF,
QQ-plot vs $\mathcal N(0, 1)$. Are they convincingly white noise?
\end{exercise}

% =========================================================
\begin{exercise}[Forecasting on the stationary scale]
\textbf{(1)} On the chosen ARMA, produce a 12-step forecast of
$w_{T+1}, \dots, w_{T+12}$ with $95\%$ CI.

\textbf{(2)} Check that the forecast returns to the unconditional
mean $\E[w] \approx 0$ as $h$ grows.

\textbf{(3)} Plot the in-sample fit and the forecast with the CI
band.
\end{exercise}

% =========================================================
\begin{exercise}[Back to the original scale]
\textbf{(1)} Recall $w_t = (1 - L)(1 - L^{12}) y_t$. Inverting one
step at a time, derive the recursion that takes
$\widehat w_{T+h}$ back to $\widehat y_{T+h}$ (telescope sum on the
seasonal and the regular difference).

\textbf{(2)} Implement the back-recursion in Python. Plot $y_t$
together with $\widehat y_{T+h}$ and the back-transformed CI.

\textbf{(3)} Take the exponential to recover sales:
\[
  \widehat{\text{Sales}}_{T+h}
  = \exp\bigl(\widehat y_{T+h} + \widehat\sigma_h^2 / 2\bigr)
\]
(conditional mean) or $\exp(\widehat y_{T+h})$ (median).
Report the predicted sales for the 12 months following December
2001.
\end{exercise}

% =========================================================
\begin{exercise}[Bonus: hold-out validation]
\textbf{(1)} Split: train $=$ Jan 1995 -- Dec 2000 ($72$ obs),
test $=$ Jan -- Dec 2001 ($12$ obs).

\textbf{(2)} Re-fit the chosen ARMA on the training set (with the
same differencing recipe), forecast the test horizon, back-transform.

\textbf{(3)} Report MAE, RMSE, MAPE on the original scale. Discuss:
does the model under- or over-estimate the December peak?
\end{exercise}

% =========================================================
\section*{Marking grid (out of 100)}

\begin{center}
\begin{tabular}{p{8cm}rl}
\toprule
\textbf{Criterion} & \textbf{Points} & \\
\midrule
Stationarisation $+$ ADF/KPSS verification        & 10 \\
ACF/PACF identification of three candidates       & 15 \\
MLE estimation $+$ AIC/BIC selection              & 20 \\
Residual diagnostics (time, ACF, QQ, LB)          & 15 \\
Forecast on the stationary scale + CI             & 10 \\
Back-transform to the original scale              & 10 \\
Bonus: hold-out validation (RMSE/MAPE)            & 10 \\
Reproducibility + report quality                  & 10 \\
\midrule
\textbf{Total}                                    & \textbf{100} \\
\bottomrule
\end{tabular}
\end{center}

\bigskip
\hfill\auteurNom\par
\hfill\auteurEmail\par

\end{document}
